% !TeX root = ../traffic_rules.tex

\section{Preliminaries}
\label{sec:preliminaries}
\subsection{Legal information}
Please note that no official translated version of the StVO exist. The translation used for this paper is based on the \textit{German Law Archive}\footnote{https://germanlawarchive.iuscomp.org/?p=1290}.
The formalized rules in this work consider the ego vehicle as the autonomous vehicle from whose perception the rules are written. All other traffic participants can either be controlled by humans or can also be autonomous vehicles.

We assume that the considered interstates have no intersections, driving directions are structurally separated so that only lanes with the same driving direction are adjacent, and have right-hand traffic. However, our formalization can also be used for left-hand traffic.
Please note that it is possible on German interstates that no speed limit exist. In the following sections, we differentiate main carriage ways, access ramps, exit ramps, and shoulder lanes as potential lane types and solid, dashed, broad solid, and broad dashed lines as possible lane markings. We consider only cars and trucks as vehicle types. %We do not consider emergencies, e.g., broken vehicle elements, and rules which restrict .


\subsection{Road Network}
We describe our road network by lanelets, which are atomic, drivable road segments, geometrically represented by a left and right bound, where the bounds are represented by polylines \cite{Bender2014}. 
The set of all lanelets in a road network is denoted by $\mathbb{L} \cup \{\bot\}$, where $\bot$ is the bottom element for the case that no lanelet exist.
Subsequently, we refer to a single lanelet as $l$, where $l \in \mathbb{L}$. A lanelet can be described by attributes, where the set of all attributes is denoted by $\mathbb{A}$. For interstates, we have $\mathbb{A} =  \{access\_ramp,$ $exit\_ramp,$ $main\_carriage\_way,$ $shoulder \}$. The function $\laneltAttribute(l)$ returns the attribute which is valid for a lanelet (cf. Fig.~\ref{fig:road_network}). Similar, the left and right bounds of a lanelet can have lanelet markings. We denote the set of possible lanelet markings as $\mathbb{LM} = \{solid, dashed, broad\_solid, broad\_dashed \}$. The functions $\mathrm{mark}_\mathrm{l}(l)$ and $\mathrm{mark}_\mathrm{r}(l)$ return the left and right lanelet marking of a lanelet. 
The functions $\laneletLeftBound(l)$, $\laneletRightBound(l)$, $\laneletBegin(l)$, and $\laneletEnd(l)$ return the left bound, right bound, left and right initial points, and left and right final points of a lanelet. The functions
\begin{equation}
\mathrm{adj}_\mathrm{l}(l) =
\begin{cases} 
l' & \text{if } l' \in \mathbb{L} \land \laneletLeftBound(l) = \laneletRightBound(l')  \\ 
\bot & \text{otherwise} 
\end{cases}	
\end{equation}
\begin{equation}
\mathrm{adj}_\mathrm{r}(l) =
\begin{cases} 
l' & \text{if } \ l' \in \mathbb{L} \land \laneletRightBound(l) = \laneletLeftBound(l')  \\
\bot & \text{otherwise} 
\end{cases}	
\end{equation}
\begin{equation}
\mathrm{succ}(l) =
\begin{cases} 
l' & \text{if } \ l' \in \mathbb{L} \land \laneletEnd(l) = \laneletBegin(l')  \\ 
& \land \, \laneltAttribute(l) = \laneltAttribute(l')   \\ 
\bot & \text{otherwise} 
\end{cases}	
\end{equation}
\begin{equation}
\mathrm{pre}(l) =
\begin{cases} 
l' & \text{if } \ l' \in \mathbb{L} \land  \laneletBegin(l) = \laneletEnd(l')  \\ 
& \land \, \laneltAttribute(l) = \laneltAttribute(l')   \\ 
\bot & \text{otherwise} 
\end{cases}	
\end{equation}
return the adjacent left, adjacent right, successor, and predecessor lanelet of $l$, respectively (cf. Fig.~\ref{fig:road_network}). The functions $\speedLimitFunction(ls)$ and  $\requiredSpeedFunction(ls)$, where $ls \subseteq \mathbb{L}$, return a set of maximum and minimum speed limits which a vehicle is allowed to drive on a provided set of lanelets. The functions $\laneletOrientation(l, s_\mathrm{l})$ and $\laneletWidth(l, s_\mathrm{l})$, where $s_\mathrm{l}$ is a longitudinal position, return the orientation of the reference path $\theta_\Gamma$ and the width of a lanelet at $s_\mathrm{l}$, respectively.
\begin{definition}[Lane]
	\label{def:lane}
	The lane originating from a lanelet is defined recursively:
	\begin{align*}
	\lane(l) = & \big( \{ \mathrm{succ}(l), l, \mathrm{pre}(l) \}  \\
	&\cup \lane(\mathrm{pre}(l)) \cup \lane(\mathrm{succ}(l)) \big) \setminus \{\bot\}
	\end{align*}
\end{definition}
\begin{figure}[!tpb]
	\def\svgwidth{0.75\textwidth}
	\centering
	\mbox{\footnotesize
		\import{figures/}{road_network.pdf_tex}
	}
	\caption{Road network defined by lanelets. The lanelets 1-6 form a lane. Lanelets 2, 4, 7, and 8 are the predecessor, successor, left adjacent, and right adjacent lanelet of lanelet 3, respectively. Lanelets 1-8 are of type $main\_carriage\_way$ and lanelet 9 is of type $access\_ramp$.}
	\label{fig:road_network}
\end{figure}
We use a curvilinear curvilinear coordinate system \cite{Hery2017} as shown in Fig.~\ref{fig:curvilinear_coordinate_system}.

\subsection{Vehicle configuration}
The state $x$ of a vehicle consists of a longitudinal state $\stateLon \in \mathbb{R}^n$ and a lateral state  $\stateLat \in \mathbb{R}^n$. The longitudinal state $x_\mathrm{lon} = \begin{bmatrix} s & v & a & j\end{bmatrix}^T$ consists of position $s$, velocity $v$, acceleration $a$, and jerk $j$.  
The lateral state $x_\mathrm{lat} = \begin{bmatrix} d & \theta \end{bmatrix}^T$ consists of the lateral distance to the reference path $\Gamma$ and orientation $\theta$ (cf. Fig.~\ref{fig:curvilinear_coordinate_system}). 
\begin{figure}[!tpb]
	\def\svgwidth{0.75\textwidth}
	\centering
	\mbox{\footnotesize
		\import{figures/}{kinematics.pdf_tex}
	}
	\caption{
		A curvilinear coordinate system aligned to the reference path $\Gamma$ with orientation $\theta_\Gamma$. The vehicle is described by the longitudinal position $s$, the lateral deviation of the reference path $d$, and the orientation $\theta$.}
	\label{fig:curvilinear_coordinate_system}
\end{figure}
The input of the vehicle is denoted by $u(t)$. The underlying vehicle model can be arbitrary as long as a transformation to our state representation is possible. The occupancy of a vehicle for a given state is denoted by $\mathcal{O}(x(t))$. The functions  $\frontPositionFunction(x(t))$, $\rearPositionFunction(x(t))$, $\rightPositionFunction(x(t))$, and $\leftPositionFunction(x(t))$ return the position of the front bumper, rear bumper, rightmost point and leftmost point of a vehicle and $\speedLimitType(x(t))$ returns the maximum allowed velocity for a vehicle type, e.g., a truck. All variables associated to the ego vehicle are denoted by the subscript $\square_\mathrm{ego}$, all variables associated to other vehicles are denote by the subscript $\square_\mathrm{o}$, and arbitrary vehicles including the ego vehicle are denoted by the subscript $\square_\mathrm{k}$ and $\square_\mathrm{p}$, where $\mathrm{o} \in \{1...N\}$, $N$ is the number of other traffic participants, and $\mathrm{k}$, $\mathrm{p} \in \{0...(N+1)\}$ represent all vehicles in the road network including the ego vehicle.
For a given initial state $x_0$, an input trajectory $u(\cdot)$, we introduce the solution of the model $\dot{x}=f(x,u)$ of the ego vehicle over time $t$ as $\xi(t;x_0,u(\cdot))$. 
We denote a braking input by $u_\mathrm{brake}(\cdot)$. The time at which a safe state is reached for $u(\cdot)$, given the system dynamics and $x_0$, is denoted by $t_s(u(\cdot))$. We consider a state to be safe on interstates for the ego vehicle if the ego vehicle is in standstill.


\subsection{Finite Linear Temporal Logic} 
We use Finite Linear Temporal Logic (FLTL) since it is interpreted over finite traces of Boolean elements.
Given is a set $\mathbb{AP}$ of atomic propositions, where each atomic proposition $\sigma_i \in \mathbb{AP}$ represents a Boolean statement. A FLTL formula $\phi$ is defined as
\begin{align*}
\phi ::= & \, \sigma_i \,|\, \neg \phi \,|\, \phi_1 \land \phi_2 \,|\, \phi_1 \lor \phi_2 \,|\,\\
& X(\phi) \,|\, \overline{X}(\phi) \,|\, \phi_1 U \phi_2 |\, G(\phi) |\, F(\phi)
\end{align*}
where the letter X stands for the weak next operator, $\overline{\text{X}}$ for the strong next, U for the until, G for the globally, and F for the future operator. We also require the Boolean operators $\neg$, $\land$, and $\lor$. Unary connectives have precedence over binary connectives, $\neg$, $X$, and $\overline{X}$ have equal precedence and $U$ has precedence over $\land$ and $\lor$. The implication $a \implies b$ is defined as $\neg a \lor b$. 

We describe the semantics of FLTL informally. For a detailed description of the semantics we refer the reader to \cite{Manna1995} and \cite{Bauer2010}. The weak next operator expresses that if a next state exists then this state has to satisfy $\phi$, whereas the strong next operator expresses that the next state has to exist and that this state has to satisfy $\phi$. The until operator says that $\phi_1$ is true at least until $\phi_2$ is true. The globally operator states that $\phi$ holds at all states and the future operator states that $\phi$ holds at some future state.

